\chapter{细胞怎样听见外界消息}

\begin{kaichang}
课堂上，老师突然点了你的名字。就在那一秒，你的心跳“咚”地快了半拍，手心开始冒汗，脸可能还微微发烫。请停下来想想这件事有多奇怪：老师的声音只是一阵空气振动，钻进你的耳朵就停了。空气振动怎么就变成了心脏加速、汗腺开工？你的耳朵和心脏之间，没有连着任何一根电线。

再看窗台上那盆含羞草。手指轻轻一碰，成排的小叶子在几秒钟内依次合拢，叶柄慢慢垂下，像睡着了一样。植物没有神经，没有肌肉，它怎么“听见”你碰了它？

先别往下看，把你的猜测写在纸上：消息从耳朵到心脏、从指尖到叶片，中间到底经过了什么？这一章结束时你会发现，这两件看似毫不相干的事，背后用的是同一套化学逻辑——一套用一行箭头就能画完、却支撑着全身几十万亿个细胞协调工作的逻辑。
\end{kaichang}

\section{现象层：消息一到，行为就变}

先不谈分子，只看你看得见、测得到的事实。生活里到处是这样的“触发—响应”：

\begin{itemize}[itemsep=2pt]
\item 从亮处走进暗房间，你的瞳孔会在一两秒内放大；用手机灯照一下眼睛，它立刻缩小。这不需要你“决定”，全自动。
\item 第 55 章见过的场景：饭后一小时测血糖，读数升高；再过两小时，它悄悄回到饭前水平。身体里有什么东西“知道”血糖高了，并下令回收。
\item 快跑两百米，心脏怦怦直跳；停下几分钟后，心跳自己慢下来。没人给心脏打电话，可它收放自如。
\item 手指划破的地方，过一会儿发红、发热、微微肿起——伤口附近的一小片区域收到了“这里有情况”的通报。
\item 含羞草被碰后叶子合拢；窗台上的花，嫩芽总是朝窗外有光的方向歪过去。植物也在响应消息，只是动作慢。
\end{itemize}

这些现象有三个共同点，值得先记下来。第一，都有明确的\textbf{触发}：光、糖、碰触。第二，响应是\textbf{有选择}的：光照瞳孔会缩，光照手心，手心什么也不会做——不是所有细胞都听同一种消息。第三，响应会\textbf{消退}：心跳会慢回来，血糖会降回来，含羞草的叶子过十几分钟会重新张开。消息不只是一声“开工”，还必须能“收工”。

想亲眼抓到一个响应，连试剂都不用。

\begin{shiyan}
\textbf{看看瞳孔这扇自动门}

材料：一面小镜子，一只手电筒（手机闪光灯即可），一间能变暗的房间。

步骤：在暗处待两分钟，让眼睛适应。然后举起镜子，用手电筒从脸颊侧面快速照向眼睛（不要正对着眼球长时间直射），同时盯住镜子里自己的瞳孔。移开灯，再观察一次。你也可以请家人配合，从侧后方观察，看得更清楚。

观察点：灯亮的一瞬间，瞳孔直径在多长时间内变化？变化有多快？移开灯后多久恢复？试试连续照三次，第三次的变化和第一次一样大吗？

安全提示：只用普通手电短时照射，严禁用激光笔，也不要直视太阳或强光源。
\end{shiyan}

\section{粒子层：受体是一把分子门铃}

镜头推进到细胞。先看“送信的一方”。

让身体各部位协调行动的化学消息，载体是一类分子，统称\textbf{信号分子}。它们的个头差异很大：神经递质乙酰胆碱只有一百多的分子量；激素胰岛素是一条 51 个氨基酸的小蛋白质（第 48 章）；而让伤口红肿的组胺，小得只有一个侧链加几个原子。个头不同、化学性格不同，但它们干的是同一件事：\textbf{被某个特定的蛋白质认出来}。

负责“认”的蛋白质叫\textbf{受体}。第 51 章说过，细胞膜上嵌满了蛋白质，其中一类就是受体：它们横跨膜的两面，细胞外那一端负责认人，细胞内那一端负责传话。认出信号分子靠的是什么？不是什么神秘的“识别功能”，就是第 48 章讲蛋白质结合时的同一套非共价作用：形状互补、电荷相吸、氢键、疏水效应。信号分子恰好嵌进受体表面一个形状相合的口袋，被一群弱作用力托住——因为弱，所以\textbf{可逆}：信号分子待一会儿又会漂走。这一点至关重要：可逆，消息才有“结束”。

真正的动作发生在结合之后：信号分子嵌进口袋，受体的整体形状会随之改变一点点——就是第 48 章说的变构。这个微小的变形，从膜外一直传到膜内那一端，于是膜内的化学环境“看见”了一个变化。请认真体会这里的因果链：\textbf{受体并没有“理解”任何消息，它只是被改变了形状}。门铃不知道来客是谁，它只是被按下去，电路随之接通。

还有一个小分支：不是每种信号都走门外的门铃。水溶性的信号（胰岛素、肾上腺素、大多数递质）穿不过磷脂双层，只能用膜受体；而个头小、疏水的信号（比如第 47 章提过的类固醇激素）能直接渗过膜，进入细胞质甚至细胞核，和细胞内部的受体结合，直接调拨基因的表达。这一类“入户送信”的信号在第 70 章讲基因开关时还会再见。

受体有多少？一个细胞表面，针对某一种信号的受体可以有\textbf{几万到几十万个}；而许多激素在血液里的浓度低到纳摩尔量级（每升 $10^{-9}$ 摩尔）就足以引起响应。做个数量级感受：$\ud{10^{-9}}{mol}$ 大约是 $6\times10^{14}$ 个分子，撒进一升水，相当于把一小撮盐撒进一个标准游泳池——而细胞照样能“听见”。灵敏度这么高，是因为接下来要讲的机制：消息进门之后会被放大。

\section{规则层：一根磷酸接力棒把消息放大}

受体把消息递进了门。门内靠什么传话？几乎所有细胞都用两套通用机制，而且常常混着用。

\textbf{第一套：磷酸化级联。}第 50 章说过，ATP 的重要本事是把磷酸基团转移给别的分子，从而改变对方的反应性。细胞里有一大类酶专门干这个：叫\textbf{蛋白激酶}，它们把 ATP 上的磷酸基团“贴”到某个靶蛋白的特定位点上。贴上磷酸基团，相当于给蛋白质别上一枚带两个负电荷的大徽章，局部电荷和形状立刻改变，这个蛋白质的活性随之打开或关闭——还是变构那一套逻辑。有贴就有拆：另一类酶叫\textbf{磷酸酶}，专职把磷酸基团拆下来，让开关复位。

关键在于，被激活的激酶本身也是酶，而一个酶分子可以一个接一个地催化很多次（第 49 章：酶干完活自己不被消耗）。于是通路可以排成接力：受体激活激酶甲，每个激酶甲激活一大批激酶乙，每个激酶乙又激活一大批激酶丙……消息像雪崩一样逐级放大，这就是\textbf{磷酸化级联}。

\textbf{第二套：第二信使。}有些受体激活后，会去启动膜上的另一种酶，让它大量制造一种小分子，比如环腺苷酸（cAMP，可以粗略理解为一个改了形状的 AMP）；或者打开钙通道，让细胞质里的钙离子浓度瞬间蹿升（细胞平时把钙离子严密关押在内质网这个“仓库”里，第 57 章）。这些小分子和离子个头小、跑得快，几毫秒内扩散到细胞各处，把消息广播给一大批下游蛋白质。因为它们是在“第一信使”（激素本身）之后、在细胞内部接力传话的，所以叫\textbf{第二信使}。

两套机制有一个共同的、常被忽略的底色：\textbf{信号的默认状态是“关”}。磷酸酶一刻不停地拆磷酸，cAMP 一刻不停地被降解，钙离子一刻不停地被泵回仓库。细胞里的“开”，是要不断花能量维持的；消息一停，一切自动归位。这就是为什么你的心跳能快也能慢回来——可逆性不是附带的好处，是设计的核心。

\begin{qiaoliang}
放大到底能放大多少？粗略估一估。一个肾上腺素分子结合一个受体；这一个被激活的受体，在它寿命内可以连续激活几十上百个中介开关蛋白（叫 G 蛋白）；每个开关蛋白启动一个制造 cAMP 的酶，这种酶每秒能造出上千个 cAMP；每个 cAMP 激活一个蛋白激酶，每个激酶又能磷酸化成百上千个靶蛋白。只取保守数字，三级接力、每级放大一百倍：

\[100\times100\times100 = 10^{6}\]

一个信号分子，换来上百万个蛋白质分子的状态改变。这就是为什么纳摩尔级的激素浓度足以命令肝细胞在几秒内放出成亿个葡萄糖分子。数字是估算，但逻辑是硬的：\textbf{放大不来自信号分子本身，来自“酶可以反复催化”这条第 49 章的老规律}。激素只按了一下门铃，屋里的扩音系统才是音量的来源。
\end{qiaoliang}

\section{符号层：把通路写成一行箭头}

现在可以把前面的话压缩成符号了。注意，下面用的是\textbf{信息流箭头}，只表示“谁把消息传给谁”，不是第 26 章那种表示物质转化的化学方程式——同一个箭头符号，两种含义，读的时候分清场合。

一条典型的信号通路：

\[\text{肾上腺素} \rightarrow \text{受体} \rightarrow \chem{cAMP} \rightarrow \text{蛋白激酶} \rightarrow \text{靶蛋白} \rightarrow \text{响应}\]

其中的“开”和“关”两步，写成物质变化就是：

\[\text{蛋白质} + \chem{ATP} \xrightarrow{\text{蛋白激酶}} \text{蛋白质--磷酸} + \chem{ADP}\]

\[\text{蛋白质--磷酸} + \chem{H_2O} \xrightarrow{\text{磷酸酶}} \text{蛋白质} + \text{磷酸}\]

再把它画成图，更直观：

\begin{center}
\begin{tikzpicture}[every node/.style={align=center}]
\node[draw, rounded corners, inner sep=4pt] (a) at (0,0) {信号分子};
\node[draw, rounded corners, inner sep=4pt] (b) at (2.8,0) {膜受体};
\node[draw, rounded corners, inner sep=4pt] (c) at (5.6,0) {第二信使\\或激酶};
\node[draw, rounded corners, inner sep=4pt] (d) at (8.6,0) {激酶级联};
\node[draw, rounded corners, inner sep=4pt] (e) at (11.4,0) {靶蛋白\\与响应};
\draw[->, thick] (a) -- (b);
\draw[->, thick] (b) -- (c);
\draw[->, thick] (c) -- (d);
\draw[->, thick] (d) -- (e);
\end{tikzpicture}
\end{center}

这张图和第 55 章那张代谢网络图是亲戚：代谢图画的是\textbf{物质}的流向，这张图画的是\textbf{信息}的流向。一个细胞同时活在两张图里，而且两张图互相咬合——信号通路的终点，常常正是某条代谢途径的开关。

\section{同一声门铃，不同的房间}

这里有一个初看自相矛盾、想通之后极其深刻的事实：\textbf{同一个信号分子，在不同细胞里会引起完全不同的响应}。

就拿出名的肾上腺素。当你被点名、受惊吓、起跑冲刺时，肾上腺把它倒进血液，流遍全身。然后：心肌细胞收到它，收缩得更快更有力；肝细胞收到它，立刻拆糖原、往血液里放葡萄糖，给肌肉送燃料；胃肠的平滑肌收到它，反而放松下来，消化放慢——紧急时刻，饭可以等会儿再消化；瞳孔开大肌收到它，瞳孔放大，多进光线。同一种分子，四种相反的命令，全都“正确”。

秘密不在信号，在听信号的细胞。第一，不同细胞表面装着\textbf{不同亚型的受体}，就像同一款门铃配了不同的室内线路。第二，就算受体相同，细胞内部的激酶组合、靶蛋白、可调用的基因程序也各不相同。信号分子只负责“按铃”，屋里放什么音乐，由住户自己的设备决定。所以“肾上腺素的作用是什么”这个问题其实没有单一答案——答案是“看在哪个细胞”。

这个原理有一个严肃的现实推论：想通过吃药只影响身体的一个部位，非常困难。治哮喘的药扩张支气管，可心脏上也有类似的受体，于是心慌手抖是常见副作用；降压药拦住了血管的信号，可能顺带影响别的器官。\textbf{药物的副作用，很大程度上就是“同一声门铃在别的房间里也响了”}。

\section{三种邮路：广播、专递和邻里传话}

信号分子投递的方式，决定了消息传多快、传多远、持续多久。身体主要用三种“邮路”：

\begin{table}[htbp]
\centering
\begin{tabular}{llll}
\toprule
邮路 & 典型信使 & 速度与范围 & 持续特点 \\
\midrule
内分泌（血液广播） & 胰岛素、肾上腺素 & 秒到分钟，可达全身 & 较持久，缓慢消退 \\
突触（点对点专递） & 乙酰胆碱等神经递质 & 毫秒级，只到下一个细胞 & 极短，毫秒级清除 \\
旁分泌（邻里传话） & 伤口处的组胺 & 秒级，只到几毫米内 & 局部、短程 \\
\bottomrule
\end{tabular}
\end{table}

\textbf{激素}走内分泌邮路：从腺体分泌进血液，随血流广播到全身，但只有装着对应受体的细胞会响应——广播的内容人人能“收到”，能“听懂”的才有反应。第 55 章里调节血糖的胰岛素和胰高血糖素，就是这条邮路上的一对老搭档。

\textbf{神经递质}走突触专递：两个细胞之间的缝隙只有约 $\ud{20}{nm}$，递质从一侧释放，扩散到这个缝隙对面，毫秒之内被接收、随即被降解或回收。范围小到只影响一个邻居，速度快到能支撑你弹钢琴。它本质上是“激素逻辑”的极速点对点版本。

\textbf{局部信号}走旁分泌邮路：手指划伤时，伤口附近的肥大细胞释放组胺，只通知周围几毫米：血管扩张一点、管壁渗漏一点，让免疫细胞和修复材料尽快赶到现场——你看到的红肿热，就是这条本地消息的效果。

植物没有神经，用的是电信号与激素的组合拳：含羞草靠离子流动产生的电信号快速传令，而朝向光源那种慢动作，靠的是生长素在茎两侧分布不均、让一边的细胞长得更快。邮路不同，信封不同，但“分子按铃—细胞响应”的逻辑一模一样。

\section{反馈、适应与噪声}

信号系统如果只负责“放大”，早就不堪重负了。它同时装备了三个“减”的机制。

\textbf{负反馈}：响应的结果反过来抑制通路本身。第 55 章在代谢网络里见过反馈抑制，信号通路里同样遍地都是：级联末端的产物常常会给开头的受体或激酶贴上“静音”标签。血糖回升后，胰岛素的分泌和作用就会被压回去。反馈让信号成为一条会自我刹车的回路，而不是一根只管踩油门的直线。身体层面的体温、血糖、血压纠偏，是这套逻辑的放大版——第 60 章专门讲它。

\textbf{适应（脱敏）}：如果同一种刺激持续存在，细胞会主动调低音量——受体被磷酸化标记、被内吞进细胞暂时“下架”，响应逐渐减弱。你走进一家花店，几分钟后就闻不到花香了，不是香味消失了，是你的嗅觉受体适应了。这个机制让细胞对“变化”敏感，而不是对“总量”敏感——对生存而言，“情况变了”通常比“情况一直是这样”更值得报警。

\textbf{噪声}：把镜头推到单个分子，信号系统其实抖得厉害。激素浓度这么低，受体与信号分子的每一次结合都是概率事件；单个细胞里某种激酶可能只有几千个分子，随机涨落不可忽略。细胞的对策和你拍照对付噪点的办法一样：\textbf{平均与阈值}——在时间上和分子数量上做平均，并设置“不到门槛不响应”的关卡。宏观世界的稳定响应，底下垫着微观世界的大量随机事件。这个画面你应该眼熟：从第 9 章的电子云、第 21 章的分子热运动一路走来，统计平均压过个体涨落，是这本书反复出现的主题。

\begin{zhengju}
“信号不进门，门里怎么知道？”——这个问题在 1950 年代难住了整个生物学界。当时已经知道，肾上腺素能命令肝细胞把糖原拆成葡萄糖，可种种迹象表明，肾上腺素似乎\textbf{并不进入细胞}。不进门的客人，怎么指挥屋里干活？

美国生化学家厄尔·萨瑟兰（Earl Sutherland）设计了一个漂亮的拆解实验。他把肝细胞打碎，用离心把“膜碎片”和“细胞质”两堆分开，然后分别测试：单独往细胞质里加肾上腺素——没反应；单独的膜碎片——也没反应。可是把\textbf{两堆重新混在一起}再加肾上腺素，糖原分解的酶立刻被激活了。

这个结果既奇怪又信息量大：肾上腺素只需要膜在场，说明它确实只作用于膜；但膜必须能\textbf{制造某种可扩散的东西}，把消息带进细胞质。萨瑟兰接着把混合液加热煮沸——蛋白质都变性失活了，可那种“传话的东西”居然还有活性。耐热的不是蛋白质，是一种小分子。1957 年前后，他和同事鉴定出这个分子：环腺苷酸，cAMP。

当然也有别的解释需要排除：会不会是肾上腺素先进了细胞、在细胞内被改造后才起作用？用带放射性标记的肾上腺素做追踪，发现它几乎只停在膜上，不进去。“第二信使”的概念就此确立：激素是第一信使，只在门外按铃；cAMP 是第二信使，在门内传话。萨瑟兰因此获得 1971 年诺贝尔生理学或医学奖。这个故事最该记住的是那个转折点：一个“分开就没活性、混起来才有”的实验，看起来像实验做砸了，实际上是一个全新概念的入口——\textbf{失败样式的数据，常常是自然在提示你分类方式错了}。
\end{zhengju}

\section{边界：通路图是一张被平均过的地图}

本章的“信号—受体—级联—响应”模型好用，但有几条边界要划清楚：

\begin{itemize}[itemsep=2pt]
\item 真实的信号系统是一张\textbf{互相串话的网络}，不是孤立的直线。一条通路的中间产物常常影响另一条通路，教科书为了讲得清楚才把它们拆成一根根独立的链条。
\item 通路图是\textbf{千万次测量的平均值}。单个细胞、单个时刻里，分子数量的涨落可能让行为偏离“标准剧本”——前面说的噪声就在这里回来。
\item 响应与浓度之间不是“全或无”的开关，而是一条逐渐上升的 S 形曲线（见本章末的挑战层）。“有没有响应”其实是“响应强到什么程度”。
\item 锁钥模型只是第一层近似。和酶一样（第 49 章），受体与信号分子结合时会互相改变形状，是“诱导契合”，不是僵硬的钥匙插锁。
\item 同一种药在不同人身上效果不同：受体数量、代谢酶活性、年龄和身体状况都有个体差异。所以涉及激素、神经递质的药物，剂量调整必须遵医嘱——本章只解释原理，不提供任何用药建议。
\end{itemize}

\begin{wuqu}
“信号分子越多，细胞响应就越强、越持久——所以这类东西补得越多越好。”这是把门铃当成了音量旋钮。反例就在“适应”机制里：长期高水平的同一种信号，会推动受体脱敏、下架，响应反而变钝。2 型糖尿病中的“胰岛素抵抗”就包含这个成分——血液里胰岛素并不少，靶细胞却对铃声越来越迟钝（它的成因复杂，受体下调只是环节之一）；长期滥用外源激素类药物，会同时压制身体自身的分泌和受体数量，一旦停药反而陷入匮乏。对信号系统而言，\textbf{剂量和节奏本身就是信息的一部分}，细胞在乎的是“消息的变化模式”，不是“分子的总吨数”。
\end{wuqu}

\section{当门铃被篡改：癌细胞怎样伪造生长命令}

正常细胞很“守纪律”：分裂增殖需要收到生长信号——一类叫生长因子的蛋白质按响对应的门铃，细胞才启动分裂程序。这是多细胞身体的秩序根基：要不要长，由集体信号说了算，不由单个细胞自己说了算。

癌细胞的问题，本质上常常是\textbf{信号系统被钻了空子}，而且花样都对应着本章讲过的零件：

\begin{itemize}[itemsep=2pt]
\item \textbf{门铃卡死在“响”的位置}：某些乳腺癌细胞把一种生长信号受体（HER2）复制了几十倍，密密麻麻插在膜上，等于门铃永远被按着，细胞不停地收到“快长”。
\item \textbf{开关卡死在“开”的位置}：有一种叫 Ras 的开关蛋白，正常工作时传一次话就自动复位；突变之后复不了位，永远停在“开”。大约三成的人类癌症带着 Ras 家族的突变——一个分子的一个小改动，把“临时消息”变成了“永久命令”。
\item \textbf{刹车被拆}：还有一类蛋白质负责在信号异常时喊停，或者勒令细胞自我了断。这些“刹车片”失效后，细胞对停止命令充耳不闻。
\end{itemize}

请注意这里的因果方向：癌细胞没有“变坏”“背叛身体”——细胞不会想。是分子开关坏了，信号通路的统计逻辑被瓦解，细胞只是在执行一套被篡改的化学程序。

坏消息里有好消息：既然知道了哪个开关卡住了，就可以设计分子去堵它。治疗慢性粒细胞白血病的伊马替尼（商品名格列卫），就是一个形状恰好能卡进那个失控激酶的口袋、强迫它静音的分子。它把这种曾经凶险的白血病变成了多数人可以长期控制的慢性病，是“理解信号通路就能设计药物”的里程碑。当然，癌细胞群体还会在药物的选择压力下演化出耐药性——一个分子开关倒下去，另一个突变站起来。这场攻防战本身就是演化在细胞群体中的上演，下一章还会从合作与演化的视角再看一次癌症。

\section{回到开场：那声点名和那株含羞草}

现在兑现承诺。老师点你名字的一秒里，发生了什么？

声波振动耳蜗里的毛细胞，细胞膜上的离子通道被机械地拉开（第 51 章的膜电位登场），电信号沿神经飞奔，一站站用神经递质接力，抵达肾上腺；肾上腺把肾上腺素倒进血液这条“广播邮路”；几秒钟内，信号流遍全身——心肌细胞的受体被按下，cAMP 在细胞内涌起，磷酸化级联层层放大，心脏跳得更快更有力；汗腺、皮肤的血管各有各的受体和线路，各自响应。你体验到的心跳、出汗、脸红，是\textbf{同一个信号分子在不同房间里按响的不同铃声}。从空气振动到心跳加速，中间没有电线的部分，全部用分子接力完成了。

含羞草呢？触碰让它叶柄基部的一群细胞（叶枕）膜上离子通道打开，电信号沿着细胞传开；接着叶枕细胞把钾离子和随之而来的水，用第 51 章的泵和渗透原理快速排到细胞外——细胞失水塌软，像撒了气的气球，叶片就垂下了。过一会儿，离子被泵回去，水跟着回流，叶子重新挺起。没有神经，没有肌肉，用的是离子、水和膨压——但“按铃—传话—响应—复位”四步，一步不少。

\section{继续深挖：会听的细胞，怎样组成会合作的身体}

这一章我们只讲了“一个细胞怎样听”。可你的身体是几十万亿个细胞的集合：同一个信号，凭什么有的细胞听、有的不听？同一套 DNA（每个细胞里都一样），凭什么长出了会收缩的肌肉细胞、会放电的神经元、会分泌激素的腺体细胞？细胞又怎样排成组织、修好伤口、在必要时让出位置？\textbf{第 59 章}就回答这些问题：细胞合作形成身体。而信号系统的“纠偏”逻辑放大到全身，体温、血糖、水盐怎样被一刻不停地调回正轨，是\textbf{第 60 章}稳态的主题。至于信号通路的最后一站——受体进了细胞核，直接拧动基因的开关——那是第三册第 70 章的重头戏。你已经站在这条链的门口了。

\begin{tiaozhan}
响应与浓度之间的定量关系，可以从一个简单假设推出来。设信号分子 \chem{L} 与受体 \chem{R} 的结合是可逆平衡：\chem{L} + \chem{R} $\rightleftharpoons$ \chem{LR}，解离常数为 $K_d$。受体被占据的比例 $\theta$ 满足

\[\theta = \frac{[\chem{L}]}{K_d + [\chem{L}]}\]

检验一下：当 $[\chem{L}] = K_d$ 时，$\theta = 1/2$，所以 $K_d$ 就是“一半受体被占据”所需的浓度，它衡量亲和力——$K_d$ 越小，受体越敏感。这条曲线在浓度取对数坐标时呈 S 形：浓度升高十倍，响应只走一小段。这解释了两个现象：为什么激素在极低浓度就能起作用、又为什么药效不会随剂量无限上涨（受体总会占满，$\theta$ 到不了 1 以上）。如果进一步考虑多个受体协同、级联放大，曲线会比上式更陡——细胞的“门槛响应”就从这里来。跳过本框不影响主线。
\end{tiaozhan}

\section*{练一练}

\begin{sikaoti}
\item 画一张信息流图：从“老师点名”到“心跳加快”，标出每一棒的载体是什么（声波、电信号，还是某种化学分子），并标出“广播邮路”和“专递邮路”各出现在哪一段。
\item 一个信号分子触发三级级联，每级放大 100 倍，最终有多少个靶蛋白分子被改变状态？如果每级只放大 10 倍呢？用这个对比说明：级联的级数对放大能力意味着什么。
\item 萨瑟兰的实验中，“膜碎片和细胞质单独都无效、混合才有效”为什么强烈提示存在一种可扩散的小分子信使？如果“膜本身就是信使”，实验结果应该会是什么样？
\item 画一条曲线：横轴是时间，纵轴是细胞响应强度，画出“持续施加同一种信号、浓度不变”时响应的变化趋势，并用受体脱敏解释曲线的形状。
\item 某个突变让一种蛋白激酶在\textbf{没有}信号时也永远处于活性状态。预测携带这个突变的细胞会怎样；如果这条通路恰好控制细胞分裂，后果可能是什么？
\item 治疗哮喘的药物作用于气道平滑肌的受体，却常引起心慌、手抖。用“同一声门铃，不同的房间”解释副作用为什么难以完全避免。
\end{sikaoti}
